\subsection{Recommendations for next steps}
Our next critical step is to construct our prototype board and conduct comprehensive testing on all modules, including RF communication, range, autonomy, telemetry data integrity, code, mechanical, g-force, and mission simulations. These recommendations are part of our testing strategy and have been carefully planned. Although there are other smaller tasks to complete, we believe these are the most critical steps to take in moving forward.

\begin{enumerate}[leftmargin=1cm, itemindent=0.25cm, noitemsep, topsep=0pt]
\item Improve and expand the PDR document, including defining every module and design element with formulas, sources, and references if time allows.
\item Conduct an After Action Review (AAR) to analyze what went well, what could be improved, what is missing, and why, and find solutions for the discovered issues.
\item Ensure all tests are carried out correctly and all sensors are correctly calibrated.
\item Conduct as many mechanical testing and design optimizations as possible.
\item Improve skills in LaTeX to effectively document the more complex project, ensuring clarity and thoroughness in technical reporting.
\item Allocate additional time for review and iterative improvements, considering the increased complexity of the project.
\item Schedule comprehensive team meetings to discuss progress, concerns, and plans for the future.
\item Regularly update the project plan to reflect current progress and any changes in priorities or objectives.
\item Seek expert feedback, particularly from specialists in muon detection, atmospheric sciences, and environmental monitoring, to validate design choices and methodologies.
\item Remain open-minded and be willing to adapt plans and strategies as necessary to achieve project goals and objectives.
\end{enumerate}
